\section{Huấn luyện lại để cứu phép đo}
\label{sec:ch08-huan-luyen-lai}

\index{ROAR}%
\index{đầu vào ngoài phân phối}%
Hooker và cộng sự đặt tên cho vấn đề mà mục trước dừng lại, và họ đặt nó bằng
ngôn ngữ của học máy chứ không bằng ngôn ngữ của XAI. Mẫu bị bỏ đi một phần đặc
trưng đến từ một phân phối khác, nên dùng nó để đánh giá là vi phạm đúng giả
thiết nền của học máy, rằng dữ liệu huấn luyện và dữ liệu đánh giá cùng một
phân phối~\autocite{mroar}.

Hệ quả của việc vi phạm ấy được các tác giả phát biểu ở dạng một câu hỏi không
trả lời được. Nếu không huấn luyện lại thì không rõ hiệu năng giảm là vì ta đã
đưa vào các dấu vết lạ nằm ngoài phân phối huấn luyện, hay vì ta thật sự đã lấy
đi thông tin~\autocite{mroar}. Hai nguyên nhân ấy cho cùng một con số, và chỉ số
không phân biệt được chúng, nên con số không đọc được.

Cách sửa mà họ đề xuất, RemOve And Retrain, can thiệp vào ô thứ ba của
hình~\ref{fig:ch08-vong-lap-do}. Bỏ $k$ đặc trưng điểm cao nhất, thay bằng một
giá trị không mang thông tin, rồi \emph{huấn luyện lại} mô hình trên dữ liệu đã
sửa, và chỉ khi đó mới đo độ chính xác~\autocite{mroar}. Lập luận cho phép huấn
luyện lại nằm ở chính chữ \enquote{không mang thông tin}: một giá trị chỉ có thể
coi là không mang thông tin nếu mô hình đã được huấn luyện để hiểu nó như
vậy~\autocite{mroar}.

Phép sửa này đắt và nó thay đổi đại lượng được đo. Mỗi mức $k$ đòi một lần huấn
luyện lại, nên chi phí của một đường cong bằng chi phí huấn luyện nhân với số
điểm trên đường cong. Quan trọng hơn, mô hình sau khi huấn luyện lại không còn
là mô hình mà ta muốn giải thích. ROAR đo xem sau khi bỏ $k$ đặc trưng ấy đi thì
\emph{một mô hình mới} còn học được bao nhiêu, tức là đo lượng thông tin nằm
trong tập đặc trưng đó đối với bài toán. Đó là một đại lượng có ý nghĩa, nhưng
nó không phải quan hệ mà định nghĩa~\ref{def:ch01-do-trung-thuc} đặt ra, vốn nói
về một hành vi của một mô hình đã cho.

Hase và cộng sự sau đó siết lập luận và nói rõ nó đi xa tới đâu. Họ gọi thứ mà
ROAR mô tả bằng đúng một cái tên, \tn{đầu vào ngoài phân phối}{out-of-distribution input},
và vì các đầu vào \tn{phản thực}{counterfactual} là như vậy, tiên nghiệm của
mô hình và cả cách khởi tạo trọng số ngẫu nhiên đều lọt được vào lời giải
thích và vào chính các chỉ số đánh giá lời giải thích~\autocite{mood}. Đường
lọt ấy, theo cách sách đọc lập luận của họ, đi như sau: trên một đầu vào mà dữ
liệu huấn luyện chưa từng bao phủ, dữ liệu không quyết định đầu ra của mô hình,
nên đầu ra do những thứ còn lại quyết định, là lớp mô hình cùng bộ điều chuẩn,
tức tiên nghiệm, và cả điểm khởi tạo ngẫu nhiên của trọng số. Một chỉ số đọc
đầu ra ấy thì đọc luôn phần do khởi tạo quyết định. Câu ấy nặng hơn câu của ROAR
một bậc. ROAR nói ta không đọc được con số; Hase và cộng sự nói con số có mang
thông tin, nhưng một phần thông tin ấy đến từ những thứ chẳng liên quan gì đến
lời giải thích lẫn mô hình đã huấn luyện.

Bài báo cũng gọi thẳng tên chỉ số mà nó nhắm vào: sufficiency của ERASER, ở
dạng chỉ giữ lại $k$ đơn vị điểm cao nhất~\autocite{mood}. Nhờ vậy lời phê phán
này không dừng ở mức một nhận xét chung về họ xóa dần, mà gắn được vào một chỉ
số cụ thể đang được dùng trong một bộ đánh giá cụ thể.

Như vậy ô thứ hai và ô thứ ba của vòng lặp đã lộ ra là những chỗ không có lựa
chọn trung lập. Hai họ tiếp theo bỏ hẳn vòng lặp, vì chúng không cần một thứ tự
trên các đặc trưng để bắt đầu, và mục sau đọc họ thứ nhất trong hai họ ấy.
